% =============================================================================
% Iter 39 -- Pillar 3: G=4 vs G=32 critical-budget detection + claim-strength audit
% Driver: scripts/group_size_iter39.py + scripts/group_size_iter39_fig.py
% Inputs: experiments/results/group_size_iter31_iso_token.tsv
%         experiments/results/group_size_iter35_pair_sweep.tsv
% Outputs: experiments/results/group_size_iter39_{t_critical,aic,claim_strength,summary}.{tsv,json}
%          figures/group_size_iter39.{pdf,png}
% =============================================================================
\subsection{Iter 39 Pillar 3 Audit: Critical-Budget $T^*$ and Claim Strength}
\label{sec:group-size-iter39}

\paragraph{Question.}
Iter 31's iso-token battery showed the Wu~2025 97.6\% retention
threshold fails for $G{=}4$ vs $G{=}32$ on hard GSM8K reasoning at every
measured budget above $T{=}1$\,M tokens; iter~35's 40-cell pair sweep
shows the claim holds in 18/40 cells overall. Iter~39 asks three
further questions that sharpen the empirical verdict:
\begin{enumerate}
  \item \textbf{At what token budget $T^*$ does retention $R$ first
        drop below the Wu~2025 97.6\% threshold?} If $T^*$ falls
        inside the canonical training range ($T \in \{1, 64\}$\,M),
        Wu's recommendation is operationally false; if $T^*$ falls
        above the range, the verdict remains that the claim holds at
        measured scales.
  \item \textbf{Is the retention decay well-described by an
        exponential-decay form $R(T) = R_\infty + (1-R_\infty) \exp(-T/\tau)$,
        or by a sigmoid?} The form choice bears on whether the
        threshold crossing is sharp (sigmoid) or smooth (exponential).
  \item \textbf{How strong is the iter~35 claim that ``18/40 cells
        pass''?} Binomial CI on the pass-fraction tells us whether the
        claim is significantly different from chance.
\end{enumerate}

\paragraph{Critical-budget detection.}
We fit $R(T) = R_\infty + (1-R_\infty) \exp(-T/\tau)$ to the four
iter~31 data points (G$=$4 vs G$=$32 at $T \in \{1, 4, 16, 64\}$\,M)
with a closed-form OLS solution for $\log(R - R_\infty)$.
Table~\ref{tab:groupsize-iter39-tstar} reports the fit.
The fitted $R_\infty = 0.7273$ and $\tau = 6.64$\,M (token-budget
scale) are tight: the 4-budget fit has only two free parameters and
the closed-form solution exactly matches the asymptotic retention
observed at $T{=}64$\,M. Solving $R(T^*) = 0.976$ gives
$T^* = 0.611$\,M tokens, with a parametric bootstrap 95\% CI of
$[0.32, 3.45]$\,M (assuming $R \sim \mathcal{N}(R_{\rm fit}(T), \sigma)$
with $\sigma = 0.0215$ measured from residuals).

\begin{table}[htbp]
\centering
\small
\begin{tabular}{lcccc}
\toprule
Pair & $R_\infty$ & $\tau$ (M) & $T^*$ (M) & 95\% CI (M) \\
\midrule
$G{=}4$ vs $G{=}32$ (iter~31, GSM8K) & 0.7273 & 6.635 & 0.611 & [0.322, 3.450] \\
\bottomrule
\end{tabular}
\caption{\textbf{Critical-budget detection for $G{=}4$ vs $G{=}32$.}
$T^*$ is the budget at which the fitted retention curve first
crosses the Wu~2025 97.6\% threshold. The fitted asymptote
$R_\infty = 0.7273$ matches the iter~31 measured minimum; the
$T^*$ point estimate of $0.611$\,M lies \emph{below} the smallest
measured budget ($T{=}1$\,M), with the bootstrap 95\% CI bracketing
the 1-3\,M range. The exponential form fits substantially better
than the sigmoid by $\Delta \mathrm{AIC} = -16.5$ (favours exp by
$e^{16.5/2} \approx 3{,}800\times$ evidence ratio), confirming the
retention decays smoothly rather than via a sharp threshold.}
\label{tab:groupsize-iter39-tstar}
\end{table}

\paragraph{Critical budgets across the iter~35 sweep.}
The same procedure applied to all 10 $(G_a, G_b)$ pairs in the
iter~35 sweep (each fit on 4 budgets) yields 10 additional $T^*$
estimates. \emph{Every} pair has $T^*$ below the maximum measured
budget ($T_{\max}{=}64$\,M); the median $T^*$ across the 10 pairs
is $\leq 1$\,M, and the worst-case pair ($G{=}4$ vs $G{=}8$) has
$T^* = 1.10$\,M [0.56, ...]. Figure~\ref{fig:group-size-iter39}(b)
ranks the 10 pairs by $T^*$ and overlays the max-measured-budget
reference. The implication: at any plausible GRPO training scale
($T \geq 4$\,M tokens), the iter~35 worst-case retention curve has
already crossed below the Wu~2025 97.6\% threshold.

\paragraph{Claim-strength audit.}
We bootstrap the iter~35 40-cell pass-fraction by treating each cell
as a Bernoulli with empirical probability $p = 18/40 = 0.45$:
95\% CI on the pass-fraction is $[0.30, 0.60]$, which \emph{straddles
0.5} ($p_{\rm boot}(\hat p < 0.5) = 0.81$). The claim that ``most
$(G_a, G_b, T)$ cells pass Wu'' is therefore \textbf{not} statistically
significant at the 5\% level on this 40-cell sample. In contrast,
the worst measured cell ($G{=}4$ vs $G{=}32$ at $T{=}64$\,M,
$R{=}0.727$) has upper 95\% CI bound $R_{\rm hi} = 0.788$, which
\emph{excludes} the Wu~2025 97.6\% threshold by a margin of $-0.188$.
This is the strong, asymmetric finding: at the worst cell, Wu is
\emph{ruled out} at 95\% confidence (no plausible retention value
brings the cell back into the equivalence band); at the population
level, the 18/40 pass rate is consistent with anywhere from 30\% to
60\% pass rate under sampling noise.

\paragraph{Model-selection sanity check.}
Across the 11 fitted retention curves, the exponential form beats the
sigmoid by AIC in \textbf{10/11} cases (mean
$\Delta\mathrm{AIC}_{\rm sig - \exp} = +4.2$). The lone exception is
$G{=}4$ vs $G{=}8$, where the sigmoid wins by $\Delta\mathrm{AIC} = -6.1$.
This is consistent with the iter~35 observation that the $G{=}4$ vs
$G{=}8$ retention curve has the slowest decay (the cells are close
enough that the retention-vs-budget curve is approximately flat with
mild curvature), while all other pairs -- especially those with the
larger $G$-gap -- exhibit monotonic exponential-style decay.

\paragraph{Reproducibility (iter 39).}
The iter~39 artifacts are produced by
\texttt{scripts/group\_size\_iter39.py} (closed-form OLS
exponential fit with parametric bootstrap; sigmoid grid-search
12$\times$12 over $(T_0, k)$ in
$[T_{\min}, T_{\max}] \times [10^5, 10^8]$; AIC with $k_{\rm exp}=2$,
$k_{\rm sig}=4$; bootstrap with $N_{\rm boot}=4000$ and seed
$20260702$). Four TSV outputs are written
(\texttt{group\_size\_iter39\_t\_critical.tsv},
\texttt{group\_size\_iter39\_aic.tsv},
\texttt{group\_size\_iter39\_claim\_strength.tsv},
\texttt{group\_size\_iter39\_summary.tsv}) and one JSON
(\texttt{group\_size\_iter39\_summary.json}) alongside the
3-panel figure (\texttt{figures/group\_size\_iter39.pdf}). Source:
\texttt{scripts/group\_size\_iter39.py} + \texttt{scripts/group\_size\_iter39\_fig.py}.

\begin{figure}[htbp]
\centering
\includegraphics[width=0.95\textwidth]{group_size_iter39.pdf}
\caption{\textbf{Iter~39 Pillar 3 critical-budget synthesis.}
\textbf{(a)} The iter~31 retention curve for $G{=}4$ vs $G{=}32$ with
the fitted exponential model
$R(T) = 0.727 + 0.273 \exp(-T/6.64\text{M})$; the Wu~2025 97.6\%
threshold (red dashed) is crossed at $T^*{=}0.61$\,M, with the
bootstrap 95\% CI band (green) bracketing the 0.32-3.45\,M range.
\textbf{(b)} $T^*$ across the 10 $(G_a, G_b)$ pairs in the iter~35
sweep; every pair's $T^*$ is below $T_{\max}{=}64$\,M, and the median
is $\leq 1$\,M. \textbf{(c)} Claim-strength audit on the 40-cell
iter~35 sweep: 18/40 cells have Wu 97.6\% in the 95\% CI (95\% CI on
the fraction itself: $[0.30, 0.60]$, straddles 50\%); 0/40 cells are
TOST-equivalent at $\epsilon{=}0.02$; the worst cell definitively
\emph{excludes} Wu (gap $-0.188$ to the threshold). Source:
\texttt{figures/group\_size\_iter39.pdf}.}
\label{fig:group-size-iter39}
\end{figure}